\chapter{Governança, Validação e Atualização}
\label{chap:governanca}

Este manual combina normas técnicas, legislação, decisões institucionais e regras de curso. Sua publicação exige governança para evitar que uma minuta técnica seja confundida com ato oficial.

\section{Instâncias a consultar}

Devem participar da validação, conforme suas competências:

\begin{itemize}
  \item Biblioteca Central da UnDF, para normalização, dados catalográficos, depósito, preservação, informações gerais do repositório e atualização normativa;
  \item Pró-Reitoria e áreas acadêmicas competentes, para regras gerais de graduação e TCC;
  \item colegiados e cursos, para compatibilidade com PPCs e regulamentos;
  \item unidade responsável pelos modelos, para publicação, atualização e suporte;
  \item instâncias de ética, para fluxos envolvendo seres humanos, animais, material biológico e dados sensíveis;
  \item área de acessibilidade, para requisitos e testes dos documentos;
  \item assessoria jurídica, quando houver dúvida sobre direitos autorais, proteção de dados, responsabilidades ou forma de aprovação.
\end{itemize}

Essas atribuições indicam participações possíveis. O ato competente deve confirmar a unidade responsável e o fluxo definitivo, sem atribuir à Biblioteca Central ou a outra unidade funções ainda não aprovadas.

\section{Etapas antes da publicação institucional}

A publicação como orientação institucional definitiva somente deve ocorrer após:

\begin{enumerate}
  \item revisão normativa;
  \item revisão da Biblioteca Central da UnDF;
  \item validação acadêmica;
  \item compatibilização com PPCs e regulamentos;
  \item sincronização com os modelos editáveis e os PDFs distribuídos;
  \item teste das validações automáticas;
  \item teste de acessibilidade;
  \item aprovação por ato formal.
\end{enumerate}

\section{Manutenção}

Cada edição deve registrar versão, data e situação. O fluxo oficial de aprovação, vigência, responsabilidades, periodicidade e arquivamento ainda precisa ser definido. Revisões extraordinárias devem ocorrer sempre que norma, legislação, regulamento ou processo institucional relevante for alterado.

\section{Conteúdo editável e elementos padronizados}

Nos modelos institucionais, o estudante deve alterar somente os campos e conteúdos destinados à autoria. Os elementos padronizados devem ser preservados conforme o modelo disponibilizado e as orientações do curso.
